% !TeX root = ../../Book1.tex
% 纲要标题：### 6.5 分类不变量与唯一性
% 纲要位置：几何学计划纲要.md，第 349 行。
% 正文状态：初稿；尚未正式审核。

% 纲要内容开始
% 
% * 定向覆盖及可定向性
% * Euler 示性数与亏格
% * 标准模型的非同胚性
% * 基本群展示的独立核验
% 
% 纲要内容结束

\section{分类不变量与唯一性}
\label{g1:sec:0605}

多边形约化说明每个紧连通曲面都能变成一个标准模型，但它尚未排除同一个曲面有两种不同的标准答案。例如，一个手柄能否经复杂的剪贴变成两个交叉帽？单看 Euler 示性数无法排除这种可能，因为两者都会使球面的示性数减少 \(2\)。我们还需要记录局部定向绕行一周后是否反转，并把边界分支数与示性数一起使用。

\subsection{定向覆盖及可定向性}
\label{g1:subsec:0605:01}

回忆\cref{g1:def:surface-orientation}：在每一点有两个局部定向，坐标之间的局部符号决定哪些选择彼此相容。可定向性要求这些局部选择能在整个曲面上连续接起来。即使曲面不可定向，我们也可以把两种选择都保留为两个不同的点，得到第五章所讨论的二叶覆盖。

\begin{theorem}[定向二叶覆盖]\label{g1:thm:orientation-cover}
	设 \(S\) 是非空连通曲面，允许有边界。令
	\[
		\widehat S=\{(x,o):x\in S,\ o\text{ 是 }x\text{ 处的一个局部定向}\},
		\qquad p(x,o)=x.
	\]
	在每个坐标邻域上，把与坐标正定向相容和与之相反的选择分别作为两张开片，则这些开片给出 \(\widehat S\) 的拓扑，使 \(p:\widehat S\rightarrow S\) 成为二叶覆盖，称为\term[dingxiangfugai]{定向覆盖}[orientation covering]。曲面 \(\widehat S\) 自然可定向，映射 \(\tau(x,o)=(x,-o)\) 是反转定向的甲板变换。并且：
	\begin{enumerate}
		\item \(S\) 可定向时，\(\widehat S\) 是两个分别同胚于 \(S\) 的连通分支；
		\item \(S\) 不可定向时，\(\widehat S\) 连通，甲板群为 \(\{1,\tau\}\)。
	\end{enumerate}
\end{theorem}
\begin{proof}
	设 \(\varphi_i:U_i\rightarrow V_i\) 为坐标。用符号 \(\varepsilon=\pm1\) 表示相对于 \(\varphi_i\) 的局部定向。若 \(\varphi_j\circ\varphi_i^{-1}\) 在点 \(x\) 的符号为 \(\sigma_{ji}(x)\)，则同一点的两种表示满足
	\[
		(x,\varepsilon)_i=(x,\sigma_{ji}(x)\varepsilon)_j.
	\]
	由\cref{g1:prop:orientation-local-sign}，\(\sigma_{ji}\) 局部常值，并满足复合的乘法规则。因此三张图重叠时识别相容，且任意两张上述开片的交集在各自坐标中开。这就给出了拓扑；每个 \(p^{-1}(U_i)\) 恰分成两张片，分别同胚于 \(U_i\)，故 \(p\) 为二叶覆盖。不同底点可用 \(S\) 的 Hausdorff 性分离，同一底点上方的两点可用两张片分离；从 \(S\) 的可数坐标基提起两份基又得第二可数性，所以总空间确为曲面。

	在 \((x,o)\) 附近，选与 \(o\) 相容的底空间坐标，将它经 \(p\) 提到所在片上。这些提升坐标之间保持定向，因而定义 \(\widehat S\) 的定向。交换 \(o\) 与 \(-o\) 则反转这一定向，并且保持投影，所以得到所述甲板变换 \(\tau\)。

	若 \(S\) 已选定全局定向 \(o_x\)，则 \(\{(x,o_x)\}\) 与 \(\{(x,-o_x)\}\) 是两个不交开闭子集，投影分别为同胚；连通性保证它们就是两个分支。反过来，若 \(\widehat S\) 不连通，由\cref{g1:prop:monodromy-components}，其分支对应二点纤维上的单值轨道。曲面连通且局部道路连通，因而道路连通，满足该命题的条件。二点集上的不传递作用只能有两个单点轨道，于是每个分支恰含每个纤维中的一个点；投影在这个分支上为双射的局部同胚，故为同胚。取其逆便得到连续的局部定向选择，说明 \(S\) 可定向。不可定向时总空间因此连通，而甲板变换由一个纤维点的像决定；二点纤维最多容许两个这样的变换，已经构造出的 \(1,\tau\) 就是全部。
\end{proof}

固定 \(x_0\) 的一个局部定向 \(o_0\)，沿环路 \(\gamma\) 提升到 \(\widehat S\)。终点或者是 \((x_0,o_0)\)，或者是 \((x_0,-o_0)\)。记相应符号为
\[
	w_S([\gamma])\in\{+1,-1\}.
\]
这称为\term[dingxiangtezheng]{定向特征}[orientation character]。第五章的同伦提升说明它只依赖环路类；按先走左边的约定拼接环路，符号相乘，故
\[
	w_S:\pi_1(S,x_0)\longrightarrow\{+1,-1\}
\]
为群同态。两种起始局部定向给出的符号相同。曲面可定向当且仅当 \(w_S\) 恒为 \(+1\)；不可定向时，定向覆盖对应子群 \(\ker w_S\)。这里曲面的小圆盘或小半圆盘可缩，故它局部道路连通且半局部单连通，确实满足\cref{g1:thm:based-covering-classification}的条件。

\ProseParagraphBreak

这个符号有直接的几何读法。在一个交叉帽的带条中绕中心走一周，横向方向反转，因而得到 \(-1\)；在一个手柄的两条标准环路上都得到 \(+1\)。边界分支的绕行也给出 \(+1\)：由\cref{g1:prop:boundary-collar}，边界附近有一个半环形领圈，它本身可定向。这个观察稍后会确定定向覆盖的边界数。

\subsection{Euler 示性数与亏格}
\label{g1:subsec:0605:02}

\cref{g1:prop:euler-triangulation-independent}已经说明，三角剖分中的数 \(V-E+F\) 不随剖分改变。为了计算标准模型，可以先使用边数少得多的多边形分解，再把它细分为三角形：分一条边会同时增加一个顶点和一条边，在一个有 \(m\) 条边的面内添中心并连到边界顶点，则顶点、边、面分别增加 \(1,m,m-1\)，都保持 \(V-E+F\)。若粘合后的边或顶点有重复，先在边上增加分点，并在面内取一圈较小的多边形，把各个角分开后逐片三角化；这些同样是上述保持交错和的细分。因此这样的有限多边形分解也可用于计算 \(\chi\)。

\begin{proposition}[标准曲面的 Euler 示性数]\label{g1:prop:surface-euler-formulas}
	对\cref{g1:def:surface-standard-models}中的标准模型，有
	\[
		\chi(\Sigma_{g,b})=2-2g-b\quad(g,b\ge0),\qquad
		\chi(N_{k,b})=2-k-b\quad(k\ge1,\ b\ge0).
	\]
	这里 \(\Sigma_{0,0}=S^2\)。沿边界封上一个圆盘使 \(\chi\) 增加 \(1\)；沿两个曲面的圆盘孔口作连通和时，
	\[
		\chi(S\mathbin{\#}T)=\chi(S)+\chi(T)-2.
	\]
\end{proposition}
\begin{proof}
	球面可取四面体边界的三角剖分，故 \(\chi(S^2)=4-6+4=2\)。当 \(g\ge1\) 时，可定向标准字给出一个顶点、\(2g\) 条边和一个面，因此 \(\chi(\Sigma_{g,0})=1-2g+1\)。不可定向标准字同样给出一个顶点、\(k\) 条边和一个面，因此 \(\chi(N_{k,0})=1-k+1\)。

	封盘时，先在边界圆周上取 \(m\ge3\) 个顶点；在加入的圆盘内取一个中心，连出 \(m\) 条边并得到 \(m\) 个三角形。已有边界的顶点和边不重复计数，所以新增加的交错和是 \(1-m+m=1\)。封 \(b\) 个盘便增加 \(b\)，给出带边界的公式。

	连通和先从两曲面各去一个开圆盘，所以交错和总共减 \(2\)。把两条孔口边界用相同数目的顶点与边细分后粘合，若边界各有 \(m\) 个顶点和 \(m\) 条边，粘合使总顶点数与总边数都减少 \(m\)，面数不变。因此粘合本身不改变交错和，得到最后的等式。
\end{proof}

\begin{proposition}[有限覆盖与 Euler 示性数]\label{g1:prop:surface-finite-cover-euler}
	设 \(S\) 是紧连通曲面，\(p:T\rightarrow S\) 为 \(d\) 叶覆盖，其中 \(1\le d<\infty\)。则 \(T\) 是紧曲面，允许不连通，并且
	\[
		\chi(T)=d\,\chi(S),
	\]
	其中不连通空间的 \(\chi\) 是各连通分支之和。
\end{proposition}
\begin{proof}
	覆盖的局部同胚性质给出曲面坐标；Hausdorff 性与第二可数性也可分别由底空间及有限张覆盖片取得。每个点选一个均匀覆盖的坐标邻域 \(U\)，再在其中选较小的邻域 \(V\)，使闭包 \(\overline V\) 紧且包含于 \(U\)。紧性给出有限个这样的 \(V\) 覆盖 \(S\)。每个 \(p^{-1}(\overline V)\) 是 \(d\) 个与 \(\overline V\) 同胚的紧集之并，有限个这样的原像覆盖 \(T\)，故 \(T\) 紧。

	对 \(S\) 的一个有限三角剖分作充分多次细分，使每个闭三角形都包含在某个均匀覆盖邻域中。这可以在有限复形的度量下使用\cref{g1:lem:lebesgue-number}，因为重心细分后三角形直径趋于零。每个顶点、开边和开三角形都恰有 \(d\) 个提升；一张三角形的各个面在提升中仍按原来的包含关系相接，给出 \(T\) 的三角剖分。于是三种数分别乘以 \(d\)，交错和也乘以 \(d\)。
\end{proof}

若 \(S=N_{k,b}\)，其定向覆盖连通且可定向，边界数为 \(2b\)：每条边界的单值恒等，所以边界的原像是两个分别一叶覆盖它的圆周。设总空间为 \(\Sigma_{h,2b}\)，则
\[
	2-2h-2b=2(2-k-b),\qquad h=k-1.
\]
所以实射影平面的定向覆盖是球面，Klein 瓶的定向覆盖是环面，Möbius 带的定向覆盖是环带。它们是同一个构造的三个低维数值例子，而非三个互不相关的特殊技巧。

\subsection{标准模型的非同胚性}
\label{g1:subsec:0605:03}

\begin{theorem}[紧连通曲面的分类与参数唯一性]\label{g1:thm:surface-classification-unique}
	每个非空紧连通曲面，允许有边界，恰属于下列同胚类型之一：
	\[
		\Sigma_{g,b}\quad(g,b\ge0),\qquad
		N_{k,b}\quad(k\ge1,\ b\ge0).
	\]
	两种类型彼此不同；在同一类型中，参数也唯一。等价地，边界分支数、可定向性和 Euler 示性数这三个量共同确定其同胚类型。
\end{theorem}
\begin{proof}
	存在性是\cref{g1:thm:surface-standard-existence}的结论。由\cref{g1:prop:boundary-invariance}，同胚必须把边界送到边界，因而保持边界分支数 \(b\)。同胚还把局部定向及其相容关系带到目标曲面，所以保持可定向性。由\cref{g1:lem:band-orientation}，标准的手柄粘合可选相容的定向，而交叉帽内有定向反转环路，故可定向模型与不可定向模型不可能同胚。最后，由三角剖分无关性，Euler 示性数也是同胚不变量。在可定向类型中，
	\[
		g=\frac{2-b-\chi(S)}2;
	\]
	在不可定向类型中，
	\[
		k=2-b-\chi(S).
	\]
	这两个公式分别唯一恢复剩下的参数。反过来，三个量相同的两曲面属于同一个标准模型，故彼此同胚。
\end{proof}

可定向模型的参数 \(g\) 称为\term[kuige]{亏格}[genus]。对不可定向模型，文献也把交叉帽数 \(k\) 称为不可定向亏格；本书在有混淆可能时直接写“交叉帽数”。这两个参数在 Euler 公式中的系数不同，不能在未说明可定向性的情况下使用同一个公式。

\begin{example}[三个不变量各自保留了什么]\label{g1:exm:surface-invariants}
	圆盘和 Möbius 带都有一条边界，但前者可定向、\(\chi=1\)，后者不可定向、\(\chi=0\)。环面和 Klein 瓶都是无边界且 \(\chi=0\)，由可定向性区分。圆盘和球面都可定向，却因边界数及 \(\chi\) 不同而不等价。可定向曲面若有一条边界，则其 Euler 示性数只能取 \(1,-1,-3,\ldots\)；例如“可定向、一条边界、\(\chi=0\)”不对应任何紧连通曲面。
\end{example}

\subsection{基本群展示的独立核验}
\label{g1:subsec:0605:04}

分类的唯一性没有使用基本群的完整结构。因此，重新从标准模型计算基本群可以核对剪贴时的生成元和关系，同时看清基本群究竟保留或遗忘了哪些几何信息。下面的交换子约定仍为 \([a,b]=aba^{-1}b^{-1}\)。

\begin{proposition}[曲面群展示]\label{g1:prop:surface-group-presentations}
	为标准曲面选定基点，并用道路把各边界环路接到基点，则可选生成元使
	\begin{align*}
		\pi_1(\Sigma_{g,b})
		&\cong\left\langle a_1,b_1,\ldots,a_g,b_g,d_1,\ldots,d_b
		\ \middle|\ [a_1,b_1]\cdots[a_g,b_g]d_1\cdots d_b=1\right\rangle,\\
		\pi_1(N_{k,b})
		&\cong\left\langle c_1,\ldots,c_k,d_1,\ldots,d_b
		\ \middle|\ c_1^2\cdots c_k^2d_1\cdots d_b=1\right\rangle.
	\end{align*}
	当 \(b=0\) 时删去全部 \(d_j\)；\(\Sigma_{0,0}=S^2\) 的群为平凡群。特别地，当 \(b\ge1\) 时，
	\[
		\pi_1(\Sigma_{g,b})\cong F_{2g+b-1},\qquad
		\pi_1(N_{k,b})\cong F_{k+b-1}.
	\]
\end{proposition}
\begin{proof}
	无边界且不为球面的情形，标准字的边界识别后是一点并的有限圆周图，其基本群由\cref{g1:prop:graph-fundamental}自由生成。再把多边形内部作为一个圆盘附着，\cref{g1:prop:attach-disk}恰把边界词的正规闭包商掉，所以分别得到交换子乘积或平方乘积的关系。球面基本群平凡见\cref{g1:prop:sphere-simply-connected}。

	对于带边界的模型，把被挖去的 \(b\) 个圆盘安排在标准多边形内部，从每条新边界到外边界选互不相交、内部在面内的弧，逐条沿弧切开。这样的弧可先把各小圆盘沿树连接，再取树的一组互不交的细带中心线；平面有孔圆盘沿这些弧切开后成为一个圆盘。若闭模型为球面，则先把球面的一张圆盘当作面，另一张当作外侧填充来做同样的构造。切口的两份方向相反，连接弧因而在附着词中成对出现；把这些弧和必要的顶点连接边选入生成树后，它们不贡献自由生成元。剩下的边是各手柄或交叉帽生成元，以及各边界圆周的 \(d_j\)。圆盘的边界依次读出标准字和 \(d_1,\ldots,d_b\)；通过选择绕行方向和连接弧次序，可以使关系具有陈述中的形式。再次应用图的基本群和附着圆盘命题，得到所写展示。若某个边界的方向或接入道路换了，字会逆转或共轭，但得到的是同一群的另一组生成元。

	当 \(b\ge1\) 时，关系中 \(d_b\) 只出现一次且指数为 \(1\)，所以它被唯一表示为其余生成元的一个词。把该词代回即可消去 \(d_b\) 与唯一的关系，留下自由群。更明确地，从其余字母的自由群到所展示群有自然满同态；反向把 \(d_b\) 送到关系所规定的逆词、其他字母保持不变，关系便成立，得到互逆同态。因此并未额外引入隐藏的关系。
\end{proof}

若只想用较容易计算的群不变量核验，可以令所有生成元交换。对任意群 \(G\)，把全部交换子生成的正规子群商掉，所得 \(G_{\mathrm{ab}}\) 称为\term[jiaohuanhua]{交换化}[abelianization]；每个到交换群的同态都唯一地经过这个商。由刚才的展示可直接得到
\[
	\pi_1(\Sigma_{g,0})_{\mathrm{ab}}\cong\mathbb Z^{2g},\qquad
	\pi_1(N_{k,0})_{\mathrm{ab}}\cong\mathbb Z^{k-1}\oplus\mathbb Z/2\mathbb Z.
\]
第一个等式因为每个交换子自动消失。第二个等式的唯一加法关系为 \(2(c_1+\cdots+c_k)=0\)：用 \(u=c_1+\cdots+c_k,c_2,\ldots,c_k\) 代替原基，逆变换为 \(c_1=u-c_2-\cdots-c_k\)，这是整数可逆换基，于是只有 \(u\) 受到 \(2u=0\) 的限制。闭可定向曲面因此没有这里的二阶扭元，而每个闭不可定向模型都有一个。

\ProseParagraphBreak

还有一个不必处理整数扭元的核验。到加法群 \(\mathbb Z/2\mathbb Z\) 的同态由生成元的像决定；在这个群中，交换子和平方都自动为零。所以对闭曲面，同态分别有 \(2^{2g}\) 个和 \(2^k\) 个，两个指数都等于 \(2-\chi\)。有边界时，唯一关系变成 \(d_1+\cdots+d_b=0\)，同态数则为 \(2^{1-\chi}\)。这说明模 \(2\) 的计数能检验 Euler 公式，却不能单独分辨可定向性。

\begin{example}[同一个自由群不足以确定曲面]\label{g1:exm:surface-group-not-classify}
	一次穿孔环面 \(\Sigma_{1,1}\)、三边界球面 \(\Sigma_{0,3}\) 和一次穿孔 Klein 瓶 \(N_{2,1}\) 的基本群都是 \(F_2\)。前两者的边界数不同，第一与第三者的可定向性不同，因此三个曲面互不同胚。若只记抽象基本群，边界环路对应的共轭类以及定向特征都可能被遗忘；这正是分类定理仍然保留边界与定向数据的原因。
\end{example}
% 纲要条目排为编号小节标题，并进入目录与 PDF 书签；说明段仍保留于纲要注释中。
